\section{Experiment Planning}\label{sec:planning}

\subsection{Subjects Selection}

The \emph{system under evaluation} is the self-hosted LLM proposer. Two subjects are selected from the same model family and training recipe, isolating the architectural contrast of interest:

\begin{itemize}[leftmargin=1.2em]
\item \texttt{Qwen/Qwen3-14B-AWQ}: dense, 14.8B parameters, 9.3\,GiB at int4;
\item \texttt{QuixiAI/Qwen3-30B-A3B-AWQ}: Mixture-of-Experts, 30.5B total with $\sim$3B active per token, 15.6\,GiB at int4.
\end{itemize}

\textbf{Deviation from the originally proposed pair.} The Qwen3.6 models named in an earlier draft cannot be served on a 20\,GB card: both exceed it \emph{on weights alone}, before any KV cache (27B AWQ = 19.05\,GiB; 35B-A3B AWQ $\approx$ 22.4\,GiB). The cause is that Qwen3.6 is multimodal and the quantizer leaves the vision tower, \texttt{lm\_head} and the linear-attention projections at 16 bits: the Hub metadata reports \texttt{I32: 25.47B, F16: 3.33B, BF16: 0.52B}, i.e.\ 3.85B unquantized parameters ($\approx$7.7\,GB) on top of 12.7\,GB of int4 weights. ``27B at int4'' does not imply 13.5\,GB. Qwen3-32B-AWQ was evaluated and rejected on arithmetic: 18.0\,GiB of weights plus $\approx$0.8\,GiB of fp8 KV cache and $\approx$1.1\,GiB of CUDA context and activations totals $\approx$19.9\,GiB against a 19.99\,GiB card.

The replacement pair fits comfortably and arguably sharpens the research question: given one 20\,GB GPU, does a 30B sparse model activating $\sim$3B parameters per token outperform a 14B dense model activating all of them? That is the decision a self-hoster actually faces. Two limitations follow and are carried into Section~\ref{sec:threats}: no single publisher ships AWQ builds of both models, so the dense arm is Qwen's official quantization and the MoE a third-party build of the same method and bit-width; and total capacity is unequal (30.5B vs.\ 14.8B), so sparsity and capacity are not fully separable. Both repositories and exact revisions are recorded in the run-table metadata. The pre-registered backup pair (IBM Granite~4.1 8B vs.\ 30B) was not needed.

The \emph{case-study workload} is the \texttt{autoresearch} protocol \cite{karpathy2026autoresearch}, a single editable training file, a fixed wall-clock budget per experiment, one scalar metric, and a keep-or-revert decision, driving the training of a compact convolutional network on CIFAR-10 at a per-experiment budget set by calibration to \textbf{45\,s} (Section~\ref{sec:execution}). The upstream repository publishes the protocol and its default inner workload but not a runnable agent loop, so the loop was re-implemented for this experiment; Section~\ref{sec:execution} explains why a purpose-built loop is required for the measurement rather than merely convenient. \texttt{autoresearch} ships with \texttt{nanochat}, a small LLM training pipeline, as its default inner workload; this study substitutes the CNN and preserves the protocol's semantics. The reason is schedule: \texttt{nanochat} costs 25--35 min per inner experiment on this hardware, which does not fit the compressed 4-week timeline, while the CNN keeps the full factorial affordable. \texttt{autoresearch} is selected for its workload characteristics (single-machine operation, fixed-length inner experiments, per-iteration keep/revert outcomes, and clean ground-truth evaluation metrics), not for the quality of the research it automates, which is out of scope. The \texttt{nanochat} workload is kept as a specified extension (Section~\ref{sec:conclusions}).

\subsection{Experimental Variables}

\begin{table}[htbp]
\caption{Independent variables (factors).}
\label{tab:factors}
\small
\begin{tabular}{@{}p{1.7cm}p{1.9cm}p{2.8cm}@{}}
\toprule
\textbf{Factor} & \textbf{Levels} & \textbf{Rationale} \\
\midrule
\texttt{proposer} & dense (14B); MoE (30B-A3B) & Architectural contrast: does sparsity buy energy without losing accuracy? \\
\texttt{patience} & greedy; patience-3 & Revert on first regression vs.\ tolerate up to 3 consecutive regressions. \\
\texttt{loop\_budget} & 10; 20 & Max inner experiments per session ($\approx$17\,min / $\approx$32\,min as executed). \\
\bottomrule
\end{tabular}
\end{table}

Table~\ref{tab:factors} lists the three factors, each at two levels. \emph{Fixed variables:} workload (CIFAR-10 CNN, fixed 45\,s training budget), intra-model configuration, namely quantization (int4, AWQ for both arms) and sampling (temperature 0.7, $\mathit{top\text{-}p}$ 0.95; see Section~\ref{sec:threats} for why temperature 0 was abandoned), reasoning mode (disabled for both arms and verified by measurement), hardware assignment (training on GPU~0, proposer on GPU~1, both RTX 4000 Ada), and pinned software versions. Fixing the intra-model configuration is a deliberate consequence of the staged design (Section~\ref{sec:design}): Stage~1 isolates the architectural effect; intra-model settings become factors only in Stage~2.

\begin{table}[htbp]
\caption{Dependent variables.}
\label{tab:depvars}
\small
\begin{tabular}{p{2.6cm}p{5.2cm}}
\toprule
\textbf{Variable} & \textbf{Operationalization} \\
\midrule
$E_{\mathit{GPU}}$ (J) & Primary executed energy outcome: sum of NVML board-energy integrals for GPU~0 and GPU~1 over the session. \\
$E_{\mathit{GPU0}}, E_{\mathit{GPU1}}$ (J) & Whole-session energy of the training and proposer devices, respectively. These are device-associated, not pure phase costs, because each board remains powered throughout. \\
$E_{\mathit{meas}}$ (J) & Sensitivity scope reconstructed from archived traces: $E_{\mathit{GPU}}+E_{\mathit{CPU,pkg}}$. The AMD trace has no separate DRAM counter, so this is not full-system energy. \\
$E/\mathrm{kept}$ (J) & Pooled across sessions in a group: $\sum E / \sum |\{\text{kept}\}|$. Zero-kept sessions contribute energy to the numerator and zero to the denominator; they are never dropped. Scope (GPU or measured components) is stated with each value. \\
$E_{\mathit{nonret}}$ (J) & Phase-aligned GPU energy of reverted, rejected, errored, or rollback-discarded iterations. Unaligned session energy is reported separately rather than silently classified. \\
\texttt{val\_acc} & Validation accuracy on a held-out CIFAR-10 split (continuous, sampled during training). \\
Test accuracy & CIFAR-10 test-set accuracy at session end; ground-truth accuracy, y-axis of the Pareto figure. \\
Wall-clock (s) & Session duration; iterations, kept, reverted, errored counts. \\
\bottomrule
\end{tabular}
\end{table}

Two accuracy-like signals are deliberately kept distinct: the agent's keep/revert decision is its \emph{claim} of improvement (a behavioral variable and the denominator of $E/\mathrm{kept}$), whereas \texttt{val\_acc} and test accuracy are \emph{ground truth} and alone enter the Pareto analysis.

\subsection{Experimental Hypotheses}

One energy-relevance hypothesis is registered per factor (all stated as two-sided nulls; $\mu$ denotes the cell mean of the named dependent variable):

\begin{description}[leftmargin=1.2em]
\item[H$_{0}^{\mathit{prop}}$:] $\mu_{E_{\mathit{GPU}}}$ and $\mu_{\mathrm{acc}}$ (test accuracy) do not differ between dense and MoE proposers. \emph{Mechanism under test:} MoE sparsity can reduce inference energy and latency; with a fixed loop budget, proposer quality changes retained yield and therefore energy productivity rather than the number of scheduled experiments.
\item[H$_{0}^{\mathit{pat}}$:] $\mu_{E_{\mathit{nonret}}}$ and $\mu_{E/\mathrm{kept}}$ do not differ between greedy and patience-3. \emph{Mechanism:} greedy reverts noise quickly (less non-retained energy) but may discard real improvements; patience tolerates regressions longer (more non-retained energy, possibly better final accuracy).
\item[H$_{0}^{\mathit{bud}}$:] $\mu_{E/\mathrm{kept}}$ does not differ between budgets 10 and 20. \emph{Mechanism:} diminishing returns near local optima should raise energy-per-kept-mutation at higher budgets; two levels bound but do not locate the knee.
\end{description}

\subsection{Experiment Design}\label{sec:design}

The design is \emph{staged}: the architecture is studied first, alone, at a fixed intra-model configuration; intra-model settings are varied only afterwards, on the architecture the first stage identifies as preferable. This ordering attributes effects cleanly (architecture first, then configuration headroom) and keeps the factorial tractable.

\textbf{Stage 1, Phase 1 (GPU factorial).} A $2 \times 2 \times 2$ full-factorial design with 3 repetitions: 8 cells, \textbf{24 sessions}, $\approx$10 GPU-hours as executed. Cell execution order is randomized and interleaved across the run table (no blocking by cell), with a \textbf{120\,s} cool-down between sessions for thermal settling; the cool-down was raised from the proposed 60\,s after calibration showed it removes the thermal trend entirely (raw and detrended standard deviations agree to four decimals). Swapping the served model between arms costs 25\,s and occurs outside the measured window. The longest session ran 39\,min. The designated baseline cell is \emph{dense $\times$ greedy $\times$ budget 10}.

\textbf{Planned CPU-served proposer phase (pending).} The proposal specified re-running four frontier-adjacent cells at two repetitions with only proposer inference moved to CPU and training unchanged on GPU~0. This phase has not been executed while its scope awaits supervisor feedback. It is therefore a planned extension, not evidence in the present report, and should not be described as a fully CPU-only loop.

\textbf{Exploratory follow-ups.} Three completed follow-ups vary reasoning mode on both architectures (Stage~2a, 12 sessions), temperature on the MoE (Stage~2b, 12 sessions), and temperature on the dense model (Stage~2c, 6 sessions). Their $n=3$ cells support manipulation checks and exploratory effect estimates, not confirmatory equivalence claims. The originally envisaged quantization and serving-stack sweep was not executed.

Across the completed campaign, 54 valid sessions contribute 660 iterations. A separate 100-iteration MoE session and retained-checkpoint replay are prepared but not yet part of the evidence.

\subsection{Data Analysis}

Analysis proceeds in four steps. (i) \emph{Descriptive statistics}: per-cell means, medians and dispersion for all dependent variables. (ii) \emph{Measurement diagnostics}: GPU energy is plotted against duration and run-table position. The ratio of arm-mean energy to arm-mean duration is used for a transparent average-power decomposition; regression slopes are reported separately. (iii) \emph{Hypothesis testing}: Shapiro--Wilk and Levene diagnostics precede the registered ART-ANOVA for GPU energy. Test accuracy is analyzed at the session level with the same main and two-way factorial terms; effect estimates and 95\,\% intervals are reported because $n=3$ per cell leaves modest accuracy power. The Stage~2a reasoning analysis uses a session-level $2\times2$ model with HC3 intervals and is explicitly exploratory. Corrected energy-per-kept estimates are ratios of pooled sums, with stratified-within-cell bootstrap intervals; zero-kept sessions remain in every numerator. (iv) \emph{Pareto analysis} (RQ2): the non-dominated set over ($E_{\mathit{GPU}}$, test accuracy) across cell means. The keep threshold is a within-session validation decision rule, not an equivalence margin for between-cell test means, so frontier membership is treated as sample-descriptive and accompanied by uncertainty.

Two choices differ from the Green Lab analysis template. We apply the Aligned Rank Transform \cite{wobbrock2011art} rather than transforming the outcome to normality: ART preserves the factorial structure and leaves the interaction terms interpretable, whereas a normalising transform changes the scale on which effects are additive. And we do not include a random effect for the session, because sessions here are independent units rather than repeated measures on the same subject; a within-subject term would misstate the design.
